\subsection{Even walks and the canonical finite algebraic hull}

For a diversity colour $p$ put $\Gamma(p)=\langle\{p\}\rangle$ and
$B(p)=\langle A_p\rangle$; these are invariant equivalence supports.
The undirected graph of $p$-edges is connected on every
$E_{\Gamma(p)}$-class, by the definition of generated equivalence.

\begin{lemma}[Even-walk classes]\label{lem:rh-evenwalk}
The $E_{B(p)}$-classes are precisely the even-$p$-walk connectivity
classes within a $p$-component. Thus
$[\Gamma(p):B(p)]\le2$. If $r\in L\cap\Gamma(p)$, then
$r\in B(p)$.
\end{lemma}
\begin{proof}
The endpoints of a two-edge $p$-walk are equal or have colour in
$A_p$, and every colour in $A_p$ occurs as the endpoints of such a
walk. Concatenating even walks therefore generates exactly
$E_{B(p)}$. A connected undirected graph has one even-walk class
unless it is bipartite, when it has exactly two classes.
In the bipartite case, every edge within the odd part of the
corresponding two-class quotient crosses its bipartition. An allowed
triangle $(r,r,r)$ would be a three-cycle of such edges, impossible.
Thus a looped colour cannot lie in the odd part
$\Gamma(p)\setminus B(p)$.
\end{proof}

\begin{theorem}[Canonical algebraic hull for real points]
\label{thm:rh-hull}
Let $A=\{a_1,\dots,a_n\}\subseteq M$ be finite and
$x\in\operatorname{acl}_M(A)\setminus A$. Put
$p_i=\col(x,a_i)$, $k=\max\{1,|\mathcal D|\}$ and let
$R_k(3)$ be the multicolour Ramsey number for a monochromatic
triangle in a complete graph with $k$ edge colours. Then
\[
 \left|\bigcap_{i=1}^n[a_i]_{E_{\Gamma(p_i)}}\right|<R_k(3).
\]
In particular, the intersection is finite. The bound is uniform
in $A$ and $x$. This statement uses no BBGTV or CFSG, only the
standing hypotheses and the stated Erd\H{o}s and standard Ramsey
theorems.
\end{theorem}
\begin{proof}
The orbit $Z$ of $x$ over $A$ is a finite realised one-point fibre
with palette $T=\bigcap_iA_{p_i}$. The set $T$ contains no looped
colour, by the following explicit extension argument.
If a diversity colour \(r\in T\) were looped,
then for EVERY integer \(m\) the finite
complete diagram on \(A\) plus \(m\) NEW
distinct points \(x_1,\ldots,x_m\), assigning
\(\col(x_j,a_i)=p_i\) for all \(i,j\)
and \(\col(x_j,x_k)=r\) for \(j\ne k\),
would have all triangles legal.
Old--old--new triangles are copied
from the realised one-point type of \(x\);
old--new--new triangles are legal because
\(r\in A_{p_i}\) for each \(i\);
new--new--new triangles are legal because
\(r\in L\). Triangle-definedness supplies
an induced realisation, and homogeneity
moves its OLD \(A\)-part back to the
original fixed real \(A\).
Every \(x_j\) then has the same complete
one-point atomic type over \(A\) as \(x\),
and hence lies in the orbit \(Z\).
The diversity \(r\)-edges make all
\(x_j\) distinct. Since \(m\) was
arbitrary this contradicts FINITENESS
of \(Z\). No SAP or abstract infinite
Ramsey theorem is used at this step.

Put $J_A=\bigcap_i\Gamma(p_i)$. If an $E_{J_A}$-class had at
least $R_k(3)$ elements, finite multicolour Ramsey would give a
monochromatic triangle on that class, in some colour
$r\in L\cap J_A$. By Lemma~\ref{lem:rh-evenwalk},
$r\in B(p_i)=\langle A_{p_i}\rangle$ for every $i$. Looped
diagonal Helly would then give a looped colour in
$\bigcap_i A_{p_i}=T$, a contradiction. Every $E_{J_A}$-class
therefore has size strictly less than $R_k(3)$. Since $x$ and
$a_i$ lie in the same $E_{\Gamma(p_i)}$-class for each $i$,
the displayed intersection is exactly the $E_{J_A}$-class of
$x$, with the asserted uniform bound.

If $A=\varnothing$, transitivity gives no algebraic real point
when $M$ is infinite. When $M$ is finite, its atomic colour graphs
are bipartite by Lemma~\ref{lem:bip}, so $M$ contains no
monochromatic triangle and $|M|<R_k(3)$ by finite Ramsey.
If $x\in A$, equality alone gives a singleton hull. These conventions
do not introduce an additional case of the theorem.
\end{proof}


\begin{corollary}[A two-power bound on canonical real hulls]
\label{cor:rh-hull-2power}
In the nontrivial case of Theorem~\ref{thm:rh-hull}, put
$J=\bigcap_{a\in A}\Gamma(\col(x,a))$.  The canonical real hull
$C=\bigcap_{a\in A}[a]_{E_{\Gamma(\col(x,a))}}$ satisfies
\[
       |C|=2^m\le 2^{|J|}\le 2^{|\mathcal D|}
       \qquad\text{for some integer }m\ge0.
\]
The same assertion holds for $A=\varnothing$ and finite $M$, with
$J=\mathcal D$ and $C=M$. This supplies an additional explicit bound
and a two-power restriction for an already finite canonical hull;
it does not replace the preceding proof of its finiteness.
\end{corollary}
\begin{proof}
The intersection $J$ of global equivalence supports is a global
equivalence support. As in Theorem~\ref{thm:rh-hull}, the displayed
hull is exactly the $E_J$-class of $x$ and is finite. Apply
Corollary~\ref{cor:rh-class-2power}. For an empty base with finite
$M$, the universal support $\mathcal D$ has $M$ as its sole class.
\end{proof}


